\begin{question}
\noindent A historian is studying the triangular layout of an ancient observation site, thought to have been used for tracking celestial events. Archaeologists have measured the distances between the three marker stones, $A$, $B$ and $C$, as shown in the diagram.

\begin{center}
\begin{tikzpicture}[scale=1, transform shape]
    \coordinate (A) at (-2,-2);
    \coordinate (B) at (4.2,-1.4);
    \coordinate (C) at (1.6,3.4);

    \draw[thick] (A) -- (B) -- (C) -- cycle;

    \node at (A) [below left, font=\small] {$A$};
    \node at (B) [below right, font=\small] {$B$};
    \node at (C) [above, font=\small] {$C$};

    \draw (A) -- (B) node[midway, below] {$14.8$ m};
    \draw (B) -- (C) node[midway, right] {$8.6$ m};
    \draw (A) -- (C) node[midway, left] {$9.3$ m};
\end{tikzpicture}
\end{center}

\noindent $AB = 14.8$ m, $BC = 8.6$ m and $AC = 9.3$ m.

\noindent Calculate the size of the largest angle in the triangle.

\vspace{5cm}

\begin{flushright}
\answerequnit{angle}{$^\circ$}{4}
\end{flushright}
\end{question}
